% Generated draft: requires mathematical and source review.
\documentclass[11pt]{article}
\usepackage[margin=25mm]{geometry}
\usepackage{fontspec}
\setmainfont{Latin Modern Roman}
\usepackage{amsmath,amssymb,mathtools,unicode-math}
\setmathfont{Latin Modern Math}
\usepackage{enumitem,needspace,xurl,hyperref}
\setlength{\parindent}{0pt}
\setlength{\parskip}{5pt}
\setlength{\emergencystretch}{3em}
\newcommand{\N}{\mathbb{N}}
\newcommand{\Z}{\mathbb{Z}}
\newcommand{\Q}{\mathbb{Q}}
\newcommand{\R}{\mathbb{R}}
\newcommand{\C}{\mathbb{C}}
\newcommand{\F}{\mathbb{F}}
\newcommand{\A}{\mathbb{A}}
\newcommand{\PP}{\mathbb{P}}
\newcommand{\E}{\mathbb{E}}
\newcommand{\eps}{\varepsilon}
\newcommand{\dd}{\,\mathrm{d}}
\begin{document}
\section*{UnsolvedMath ID: 800479}
\textbf{Source problem code:} catalogue entry
\subsection{Definitions and mathematical statement}
In the usual Euclidean interpretation of the source, let \(D\subset\R^n\) be a convex domain with smooth boundary, and let \(x\mapsto L_x\) be a smooth field of one-dimensional linear subspaces transverse to the boundary:
\[
\R^n=T_x\partial D\oplus L_x\qquad(x\in\partial D).
\]
The source does not specify the ambient dimension, boundedness, strict convexity, or the precise differentiability requirements; the notation here describes the usual setup for \(n\geq2\) without resolving those qualifications.

For \(v\in\R^n\), write uniquely \(v=v_T+v_L\), with \(v_T\in T_x\partial D\) and \(v_L\in L_x\). The projective reflection at \(x\) is
\[
R_x(v)=v_T-v_L.
\]
Thus the tangential component in this direct-sum decomposition is preserved and the transverse component is reversed. These components need not be orthogonal projections.

A boundary model of the phase space is
\[
\mathcal M=\{(x,[v]_+):x\in\partial D,\ v\text{ points strictly into }D\},
\qquad [v]_+=\{cv:c>0\}.
\]
Where the ray from \(x\) in direction \(v\) has a next transverse boundary intersection \(y\), define the billiard map by
\[
B(x,[v]_+)=(y,[R_y(v)]_+).
\]
Write \(\mathcal U\subseteq\mathcal M\) for its domain of smooth definition. Using this natural domain avoids assuming, without source support, that every trajectory has a next collision.

\textbf{Problem, attributed to Serge Tabachnikov.} Characterize the projective billiard tables \((D,L)\) that admit a smooth symplectic form \(\omega\) on their phase space invariant under the billiard map. In the boundary model above, this means
\[
\omega\in\Omega^2(\mathcal M),\qquad
\dd\omega=0,\qquad
\ker\omega_z=\{0\}\quad(z\in\mathcal M),\qquad
B^*\omega=\omega|_{\mathcal U}.
\]
The source does not specify whether additional restrictions on the form or its behavior near grazing trajectories are intended.

The related question is to characterize when the projective reflection law has a variational description. The source does not specify the admissible action functionals, variables, regularity, locality, or equivalence notion for such a description. Consequently, this second question remains a request for a characterization whose precise variational framework requires clarification; no particular generating function or length functional is imposed here.

\textbf{Supplied background and qualifications.} The source attributes the definition of projective billiards to [Ta97] and describes the following examples and related result:
\begin{itemize}
\item In [Ta97-1], the table is bounded by a sphere and its transverse field satisfies a certain exactness condition. The invariant symplectic form on oriented lines intersecting the sphere is induced by the hyperbolic metric inside the sphere, in the projective, or Beltrami--Cayley--Klein, model. The exactness condition is not supplied, so this example is not a complete sufficient criterion as stated here.
\item In dimension \(2\), the source describes a billiard curve that is an arbitrary oval with all transverse-field lines concurrent. It attributes an invariant area form, obtained from the standard planar area form by polar duality, to [Ta97]. The source does not specify the precise oval convention or duality construction; the standing transversality requirement still applies.
\item The source reports that [Glu] and [GM] study when a Minkowski billiard is projective, states that such a billiard has a variational description and an invariant symplectic form, and gives the answer as a conventional billiard in some Euclidean metric. The definitions and hypotheses needed to formulate that reported characterization precisely are absent from the supplied record. This is retained as attributed background, not as an independently verified theorem or a resolution of the general questions.
\end{itemize}
\subsection{Short English statement}
A projective billiard reflects a direction by preserving its component tangent to the boundary and reversing its component along a prescribed transverse line. Determine which such billiards preserve a symplectic form on phase space, and which admit a variational description. The source leaves the variational framework and several geometric hypotheses unspecified. Its registry status is supplied as \emph{open}; this draft does not independently verify that status.
\subsection{Sources}
\begin{itemize}
\item Local UnsolvedMath registry, AMR Open Problem Lists: permanent ID \(800479\), exact source problem code \texttt{catalogue entry}; title: \emph{Billiards and Quantitative Symplectic Geometry --- Serge Tabachnikov --- Problem 1}. The record lists Albach, Barraud et al. as proposers, gives the year as 2026, and labels the problem open.
\item Supplied source list: Albach, Barraud et al., \emph{Open Problems in Billiards and Quantitative Symplectic Geometry} (2026), Serge Tabachnikov --- Problem 1. Contributor: Serge Tabachnikov. Supplied URL: \url{https://arxiv.org/abs/2602.12896}. The record reports access on 2026-10-16 and extraction from arXiv LaTeX. These are supplied provenance details, not a claim of independent source access or verification.
\item [Ta97]: supplied reference label for the definition of projective billiards and the planar concurrent-field example. Full bibliographic details and a reference URL are not supplied.
\item [Ta97-1]: supplied reference label for the spherical-table example with an exactness condition. Full bibliographic details, a reference URL, and the exactness condition are not supplied.
\item [Glu] and [GM]: supplied reference labels for the reported characterization of projective Minkowski billiards. Full bibliographic details, reference URLs, and precise theorem hypotheses are not supplied. The description of these papers as recent belongs to the supplied source.
\end{itemize}
\end{document}
